\documentclass[11pt]{article}
\usepackage[margin=1in]{geometry}
\usepackage{booktabs}
\usepackage[hidelinks]{hyperref}
\setlength{\parindent}{0pt}
\setlength{\parskip}{6pt}
\newcommand{\code}[1]{\texttt{#1}}

\title{UCLA TA1 --- AIQ Phase 1 Submission Summary}
\author{UCLA TA1 Team (VeriStressGT)}
\date{September 29, 2026}

\begin{document}
\sloppy
\maketitle

\paragraph{Submission.}
Evaluation card \code{cards/thrust1\_soundness.yaml}, branch \code{evaluation-card} of
\url{https://github.com/dtroxell19/VeriStressGT}. Setup and
run instructions are in \code{AIQ\_README.md} in the repository.

\section{Approach}
A neural-network verifier either proves that a network is robust, meaning no small perturbation of
an input can change its output (UNSAT), or finds a perturbed input that does (a counterexample,
SAT). Verifiers are used to certify networks in safety-critical settings, so a verifier bug that
wrongly reports ``robust'' gives false assurance. Catching such bugs requires knowing the right
answer, which is what VeriStressGT provides.

VeriStressGT tests neural-network verifiers on networks whose robustness answer is known by
construction, so a wrong verdict is provably wrong. Existing benchmarks have no
ground truth and fall back on majority vote between verifiers.
\begin{itemize}
  \item \textbf{Benchmark:} 50 instances, pre-built in the repository. 24 are UNSAT (i.e. robust) by analytic
  construction. 26 are SAT (non-robust), each with a counterexample checked in float64 and onnxruntime float32.
  \item \textbf{Verifier pool:} 2 sound in-house controls; 5 real verifiers
  ($\alpha$-$\beta$-CROWN, PyRAT, nnenum, NeuralSAT, Marabou); and 42 deliberately buggy
  verifiers: 12 mutants of our reference verifier plus 6 output-level faults on each real verifier.
  \item \textbf{Scoring:} every definitive verdict is judged against the known label. The same
  verdicts are also judged by majority vote, as a baseline.
\end{itemize}

\section{Expected Results and BAA Alignment}
Every verifier in the pool gives a verdict (SAT or UNSAT) on every instance, or times out. Because
each instance's correct answer is known, the card checks each verdict and reports two metrics:
\begin{itemize}
  \item \textbf{Scoring accuracy} (\code{gt\_scoring\_accuracy}, the BAA metric, target
  $\geq 95\%$): of all SAT/UNSAT verdicts, the fraction that the labels judge correctly, i.e.\
  a correct verdict is accepted and a wrong one is rejected.
  \item \textbf{Buggy verifiers caught} (\code{gt\_detection\_rate}): of the planted buggy
  verifiers whose bug changed at least one verdict (``fired''), the fraction shown to be wrong on
  some instance.
\end{itemize}
Both are also computed with majority-vote labels instead of known labels, as the baseline that
existing benchmarks rely on.

The run should end with \code{RESULT: VERIFIED}, with both metrics at 1.0 in the top-level
\code{verdict.json}. The card passes only if the known labels score 100\% of verdicts correctly,
flag no sound control, and catch at least 75\% of the bugs that fired. Our reference run took
32 minutes on a MacBook Pro (Apple Silicon, 8\,GB, CPU only) with MAGNET v0.1.0. The card has been
run on macOS.

\textbf{Scoring accuracy} (2{,}025 SAT/UNSAT verdicts). A wrongful accusation calls a correct
verdict wrong; a wrongful acquittal calls a wrong verdict correct; a tie is a majority vote with
no winner.
\begin{center}\small
\begin{tabular}{@{}lrrrr@{}}
\toprule
Labels used for scoring & Correct & Wrongful accusations & Wrongful acquittals & Ties \\
\midrule
Known labels (VeriStressGT) & 2025/2025 (100\%) & 0 & 0 & 0 \\
Majority vote, one buggy verifier & 2015/2025 (99.5\%) & 0 & 0 & 10 \\
Majority vote, full pool & 1959/2025 (96.7\%) & 19 & 25 & 22 \\
\bottomrule
\end{tabular}
\end{center}

\textbf{Buggy verifiers caught} (42 planted, 39 fired)
\begin{center}\small
\begin{tabular}{@{}lll@{}}
\toprule
Labels used for scoring & Caught & Sound verifiers wrongly flagged \\
\midrule
Known labels (VeriStressGT) & 39/39 (100\%) & none \\
Majority vote, one buggy verifier & 37/39 (95\%) & none \\
Majority vote, full pool & 39/39 (100\%) & 3: \code{reference}, \code{ibp\_only}, real NeuralSAT \\
\bottomrule
\end{tabular}
\end{center}

\textbf{Models and datasets.} The instances span many distinct networks in three families (MLPs,
CNNs, attention models), well past the BAA minimum of two. The inputs are synthetic and part of
each network's construction.

What can differ on other hardware, without changing the verdict or the metrics:
\begin{itemize}
  \item \textbf{Verdict counts:} timeouts depend on machine speed, so totals (2{,}025 verdicts,
  39 fired bugs) can shift slightly. Ground-truth accuracy stays 100\% because the labels do not
  depend on timing.
  \item \textbf{Marabou:} Marabou 2.0.0 is nondeterministic. We saw it return false SAT on
  \code{meap\_02}, \code{meap\_03} and \code{milp\_s2\_a}, and not every run hits all three. Each is
  a real Marabou soundness bug that ground truth correctly flags.
\end{itemize}

\section{Scalability}
TA2 scales the card with one parameter, \code{scale: N}, in the card's \code{algo\_params}. The
runner builds a benchmark $N$ times larger with fresh seeds (\code{thrust1\_bench\_xN/}) and runs
the same verifier pool on it.
\begin{itemize}
  \item \textbf{What grows:} $N$ re-seeded copies of every base network, each keeping its analytic
  UNSAT certificate, plus $N$ times the random and attack-hard CNN SAT instances.
  \item \textbf{Why results stay consistent:} every label is re-derived from construction and every
  SAT witness is re-certified, so expected ground-truth accuracy is 100\% at any $N$.
  \item \textbf{Tested:} at \code{scale: 2}, 98 instances (48 UNSAT, 50 SAT); all 50 witnesses
  re-certified, ground-truth scoring 173/173, no control flagged (sound controls only, 9 minutes,
  free size-limited Gurobi license).
  \item \textbf{Needs:} \code{gurobipy} (installed by the setup script) to build larger benchmarks.
\end{itemize}

\section{Impact and Comparison with Other Benchmarks}
Existing benchmarks cannot catch most of these bugs. We ran the same verifier pool on two VNN-COMP
benchmarks, \code{mnist\_fc} (30 of its 90 instances) and \code{oval21} (all 30).

A bug is \emph{caught by a benchmark's own labels} when one of that benchmark's labels shows a
verdict to be wrong. VeriStressGT has a known label for every instance. An external benchmark only
has labels an outside party can establish: SAT where an attack finds a counterexample that
checks out numerically. Robustness (UNSAT) cannot be certified from outside, so every other
instance is unlabelled, and a wrong verdict there goes unnoticed.

\begin{center}\small
\begin{tabular}{@{}lrrr@{}}
\toprule
Benchmark & Bugs that fire & Caught by own labels & Verdicts that can be checked \\
\midrule
VeriStressGT & 39/42 & 39/42 & 100\% \\
VNN-COMP \code{mnist\_fc} (30/90) & 19/42 & 6/42 & 21\% \\
VNN-COMP \code{oval21} & 19/42 & 2/42 & 7\% \\
\bottomrule
\end{tabular}
\end{center}

\begin{itemize}
  \item \textbf{Bugs that never fire:} 19 of the 42 bugs never change a verdict on either external
  benchmark, so no labelling could catch them there. VeriStressGT exposes 16 of these 19.
  \item \textbf{Real bugs found:} unmodified Marabou 2.0.0 returns false SAT on the provably robust
  VeriStressGT instances \code{meap\_02}, \code{meap\_03} and \code{milp\_s2\_a}, and false UNSAT
  on \code{oval21} instance \code{vb\_oval21\_0014}, which a certified counterexample shows is SAT.
  \item \textbf{Subset:} more \code{mnist\_fc} instances can only make more bugs fire, so the subset
  may understate full \code{mnist\_fc}; \code{scripts/run\_external\_comparison.sh} runs all 90.
  \item \textbf{Thrust 2} (separate self-evaluation card, not the submission): per-verifier timeout
  prediction AUC is 0.77--0.92 on held-out data and 0.87--0.96 on fresh networks, all above the 0.7
  target.
\end{itemize}

\end{document}
